
\documentclass{article}
\usepackage{colortbl}
\usepackage{makecell}
\usepackage{multirow}
\usepackage{supertabular}

\begin{document}

\newcounter{utterance}

\centering \large Interaction Transcript for game `cladder', experiment `full\_v1.5\_default', episode 4422 with Qwen3.5{-}9B.
\vspace{24pt}

{ \footnotesize  \setcounter{utterance}{1}
\setlength{\tabcolsep}{0pt}
\begin{supertabular}{c@{$\;$}|p{.15\linewidth}@{}p{.15\linewidth}p{.15\linewidth}p{.15\linewidth}p{.15\linewidth}p{.15\linewidth}}
    \# & \multicolumn{2}{c}{Player} && \multicolumn{2}{c}{Game Master} \\
    \hline

    \theutterance \stepcounter{utterance}  
    & & & \multicolumn{4}{p{0.6\linewidth}}{
        \cellcolor[rgb]{0.9,0.9,0.9}{
            \makecell[{{p{\linewidth}}}]{
                \texttt{\tiny{[P1$\langle$GM]}}
                \texttt{You are an expert in causal inference. The following question is not a typical commonsense query, but rather a meticulously designed question created by a professor specializing in causal inference, intended to assess the students' mastery of the course content.} \\
\texttt{Imagine a self{-}contained, hypothetical world with only the following conditions, and without any unmentioned factors or causal relationships: Gender has a direct effect on smoking and peanut allergy. Smoking has a direct effect on tar deposit. Tar deposit has a direct effect on peanut allergy. Gender is unobserved. For nonsmokers, the probability of high tar deposit is 18\%. For smokers, the probability of high tar deposit is 64\%. For nonsmokers and with no tar deposit, the probability of being allergic to peanuts is 46\%. For nonsmokers and with high tar deposit, the probability of being allergic to peanuts is 77\%. For smokers and with no tar deposit, the probability of being allergic to peanuts is 17\%. For smokers and with high tar deposit, the probability of being allergic to peanuts is 43\%. The overall probability of smoking is 49\%. Will smoking decrease the chance of being allergic to peanuts?} \\
\texttt{Start your answer with 'Yes' or 'No', followed by additional reasoning or evidence to support your explanation.} \\
            }
        }
    }
     \\ \\

    \theutterance \stepcounter{utterance}  
    & \multicolumn{4}{p{0.6\linewidth}}{
        \cellcolor[rgb]{0.561,0.737,0.561}{
            \makecell[{{p{\linewidth}}}]{
                \texttt{\tiny{[P1$\rangle$GM]}}
                \texttt{No} \\
\\ 
\texttt{**Reasoning:**} \\
\\ 
\texttt{To determine if smoking decreases the chance of being allergic to peanuts, we must calculate the causal effect of smoking on peanut allergy. In this causal graph, **Gender** is a confounder because it has direct arrows to both **Smoking** and **Peanut Allergy**, and it is unobserved. Therefore, we cannot simply compare the observed probabilities of allergy between smokers and nonsmokers (which would be biased by confounding). Instead, we must estimate the **Average Causal Effect (ACE)** by conditioning on the mediator (**Tar Deposit**) and the confounder (**Gender**), or more simply, by calculating the effect of smoking on tar deposit and then propagating that through the tar deposit to the allergy outcome, while accounting for the fact that Gender affects both the exposure and the outcome directly.} \\
\\ 
\texttt{However, a more rigorous approach in this specific DAG structure (\$Gender \textbackslash{}to Smoking\$, \$Gender \textbackslash{}to Allergy\$, \$Smoking \textbackslash{}to Tar\$, \$Tar \textbackslash{}to Allergy\$) is to calculate the counterfactual probabilities. Since Gender is unobserved, we cannot condition on it directly. But we can analyze the mechanism:} \\
\\ 
\texttt{1.  **Effect of Smoking on Tar Deposit:**} \\
\texttt{    *   \$P(Tar=High | Do(Smoking=Yes)) = 0.64\$} \\
\texttt{    *   \$P(Tar=High | Do(Smoking=No)) = 0.18\$} \\
\texttt{    *   Smoking *increases* the probability of high tar deposit by \$0.64 {-} 0.18 = 0.46\$.} \\
\\ 
\texttt{2.  **Effect of Tar Deposit on Peanut Allergy:**} \\
\texttt{    *   For Nonsmokers: \$P(Allergy | Tar=High) = 0.77\$, \$P(Allergy | Tar=Low) = 0.46\$. (Tar increases allergy risk here).} \\
\texttt{    *   For Smokers: \$P(Allergy | Tar=High) = 0.43\$, \$P(Allergy | Tar=Low) = 0.17\$. (Tar increases allergy risk here).} \\
\texttt{    *   In both groups, having high tar deposit increases the probability of peanut allergy compared to having no tar deposit.} \\
\\ 
\texttt{3.  **The Causal Path:**} \\
\texttt{    The causal path from Smoking to Peanut Allergy is \$Smoking \textbackslash{}to Tar \textbackslash{}to Allergy\$.} \\
\texttt{    Since smoking increases the probability of high tar deposit, and high tar deposit increases the probability of peanut allergy (regardless of smoking status), the indirect causal effect of smoking on peanut allergy via tar deposit is **positive** (it increases the risk).} \\
\\ 
\texttt{4.  **The Confounding Path:**} \\
\texttt{    There is a backdoor path: \$Smoking \textbackslash{}leftarrow Gender \textbackslash{}to Allergy\$.} \\
\texttt{    The question asks if smoking *decreases} \\
            }
        }
    }
    & & \\ \\

    \theutterance \stepcounter{utterance}  
    & & & \multicolumn{2}{p{0.3\linewidth}}{
        \cellcolor[rgb]{0.9,0.9,0.9}{
            \makecell[{{p{\linewidth}}}]{
                \texttt{\tiny{[GM$|$GM]}}
                \texttt{No} \\
            }
        }
    }
    & & \\ \\

    \theutterance \stepcounter{utterance}  
    & & & \multicolumn{2}{p{0.3\linewidth}}{
        \cellcolor[rgb]{0.9,0.9,0.9}{
            \makecell[{{p{\linewidth}}}]{
                \texttt{\tiny{[GM$|$GM]}}
                \texttt{no} \\
            }
        }
    }
    & & \\ \\

    \theutterance \stepcounter{utterance}  
    & & & \multicolumn{2}{p{0.3\linewidth}}{
        \cellcolor[rgb]{0.9,0.9,0.9}{
            \makecell[{{p{\linewidth}}}]{
                \texttt{\tiny{[GM$|$GM]}}
                \texttt{game\_result = WIN} \\
            }
        }
    }
    & & \\ \\

\end{supertabular}
}

\end{document}
